\documentclass[12pt]{article}
\usepackage{amsmath,amsfonts,amssymb}
\usepackage{geometry}% 1-inch margins
\geometry{letterpaper, margin=1in}% 1-inch margins

\usepackage{titlesec}  % Allows customization of section/subsection titles
\titleformat{\section}[block]{\normalfont\Large\bfseries}{\thesection}{1em}{}
\titleformat{\subsection}[block]{\normalfont\large\bfseries}{\thesubsection}{1em}{}

\begin{document}
% Horizontal line
\hrule
\vspace{0.2cm}

% Centered title
\begin{center}
    \LARGE \textbf{EMCH 290} \\
    \Large Mechanical Engineering Thermodynamics \\
    \large HW \#1 \textbf{SOLUTIONS}
\end{center}

% Horizontal line
\hrule
\vspace{0.2cm}

% Author and points possible
\begin{center}
    \textbf{Author:} Jacob Vaught \\
    \vspace{0.2cm}
    \textbf{Semester:} Fall 2023 \\
    \vspace{0.2cm}
    \textbf{Points Possible: 100 [2\% of FINAL Class Grade]}
\end{center}

% Another horizontal line
\hrule
\vspace{0.5cm}

%Question 1
\newpage
\section*{Question 1 [16 points]}
\subsection*{Points Breakdown:}
\begin{itemize}
    \item 1 L to \(in.^3\): 2 points
    \item 650 J to Btu: 2 points
    \item 0.135 kW to \(ft \cdot lbf/s\): 2 points
    \item 378 \(g/s\) to \(lb/min\): 2 points
    \item 304 kPa to \(lbf/in.^2\): 2 points
    \item 55 \(m^3/h\) to \(ft^3/s\): 2 points
    \item 50 \(km/h\) to \(ft/s\): 2 points
    \item 8896 N to ton: 2 points
\end{itemize}

\subsection*{Solutions:}
\begin{enumerate}
    \item \(1\, L \times 61.0237\, \frac{in^3}{L} = 61.0237\, in^3\) 
    \item \(650\, J \times 9.47 \times 10^{-4}\, \frac{Btu}{J} = 0.61555\, Btu\)
    \item \(0.135\, kW \times 737.56\, \frac{ft \cdot lbf/s}{kW} = 99.57\, ft \cdot lbf/s\)
    \item \(378\, \frac{g}{s} \times 0.1322\, \frac{lb/min}{g/s} = 49.992\, lb/min\)
    \item \(304\, kPa \times 0.145\, \frac{lbf/in.^2}{kPa} = 44.08\, lbf/in.^2\)
    \item \(55\, \frac{m^3/h}{3600\, s/h} \times 35.314\, \frac{ft^3}{m^3} = 0.5397\, ft^3/s\)
    \item \(50\, \frac{km/h}{3600\, s/h} \times 3280.84\, \frac{ft}{km} = 45.565\, ft/s\)
    \item \(8896\, N \times 0.0002248\, \frac{ton}{N} = 2\, ton\)
\end{enumerate}


\newpage
\section*{Question 2 [16 points]}
\subsection*{Points Breakdown:}
\begin{itemize}
    \item \(122\, in^3\) to L: 2 points
    \item \(778.17\, ft \cdot lbf\) to kJ: 2 points
    \item \(100\, hp\) to kW: 2 points
    \item \(1000\, \frac{lb}{h}\) to \( \frac{kg}{s}\): 2 points
    \item \(29.392\, \frac{lbf}{in^2}\) to bar: 2 points
    \item \(2500\, \frac{ft^3}{min}\) to \( \frac{m^3}{s}\): 2 points
    \item \(75\, \frac{mile}{h}\) to \( \frac{km}{h}\): 2 points
    \item \(1\, ton\) to N: 2 points
\end{itemize}

\subsection*{Solutions:}
\begin{enumerate}
    \item \(122\, in^3 \times 0.0163871\, \frac{L}{in^3} = 1.99783\, L\)
    \item \(778.17\, ft \cdot lbf \times 1.35582\, \frac{kJ}{ft \cdot lbf} = 1.05519\, kJ\)
    \item \(100\, hp \times 0.7457\, \frac{kW}{hp} = 74.57\, kW\)
    \item \(1000\, \frac{lb}{h} \times 1.25998 \times 10^{-4}\, \frac{kg}{s} = 0.126\, \frac{kg}{s}\)
    \item \(29.392\, \frac{lbf}{in^2} \times 0.0689476\, \frac{bar}{\frac{lbf}{in^2}} = 2.02439\, bar\)
    \item \(2500\, \frac{ft^3}{min} \times 4.71947 \times 10^{-4}\, \frac{m^3}{s} = 1.17987\, \frac{m^3}{s}\)
    \item \(75\, \frac{mile}{h} \times 1.60934\, \frac{km}{mile} = 120.7\, \frac{km}{h}\)
    \item \(1\, ton \times 2000\, \frac{lbf}{ton} \times 4.44822\, \frac{N}{lbf} = 8896.44\, N\)
\end{enumerate}

\newpage
\section*{Question 3 [4 points]}
\subsection*{Points Breakdown:}
\begin{itemize}
    \item Calculation of the weight in \(kN\): 4 points
\end{itemize}

\subsection*{Solutions:}
\begin{enumerate}
    \item The weight of the object is found using the formula \( F = m \times a \). With \( m = 10,000\, kg \) and \( a = 9.81\, \frac{m}{s^2} \), the weight of the object in \(N\) is calculated as \( F = 10,000\, kg \times 9.81\, \frac{m}{s^2} = 98,100\, N \).
    \item To convert this weight into \(kN\), the conversion factor \(1\, kN = 1000\, N\) is used. Therefore, \( 98,100\, N \times \frac{1\, kN}{1000\, N} = 98.1\, kN \) is obtained.
\end{enumerate}

\newpage
\section*{Question 4 [5 points]}
\subsection*{Points Breakdown:}
\begin{itemize}
    \item Calculation of the mass of the satellite: 2 points
    \item Calculation of the weight of the satellite in space: 3 points
\end{itemize}

\subsection*{Solutions:}
\begin{enumerate}
    \item The mass of the satellite is determined by rearranging the weight formula \( W = g \times m \) to solve for \( m \), as \( m = \frac{W}{g} \). With \( W = 4400\, N \) and \( g = 9.8\, \frac{m}{s^2} \), the mass of the satellite is found to be \( m = \frac{4400}{9.8} = 448.98\, kg \).
    \item The weight of the satellite in space is calculated using the formula \( W = g \times m \). With the given acceleration due to gravity in space as \( g = 0.224\, \frac{m}{s^2} \) and the mass \( m = 448.98\, kg \), the weight in space is found to be \( W = 0.224 \times 448.98 = 100.57\, N \).
\end{enumerate}

\newpage
\section*{Question 5 [6 points]}
\subsection*{Points Breakdown:}
\begin{itemize}
    \item Calculation of the weight or force exerted by the piston on the gas: 2 points
    \item Calculation of the force applied by the atmosphere: 2 points
    \item Calculation of the pressure needed for the piston to rise: 2 points
\end{itemize}

\subsection*{Solutions:}
\begin{enumerate}
    \item The force exerted by the piston on the gas is calculated using the formula \( F = m \times g \). With \( m = 50 \, \text{kg} \) and \( g = 9.81 \, \frac{\text{m}}{\text{s}^2} \), the force is \( F = 50 \times 9.81 = 490.5 \, \text{N} \).
    \item The force applied by the atmosphere is calculated using \( F = P \times A \). With a given atmospheric pressure of \( 1 \, \text{bar} = 1 \times 10^5 \, \text{Pa} \) and the area \( A = 0.01 \, \text{m}^2 \), the force is \( F = 1 \times 10^5 \times 0.01 = 1000 \, \text{N} \).
    \item The pressure needed for the piston to rise is calculated using \( F = P \times A \). The total force is the sum of the forces due to the piston and the atmosphere, \( F = 1000 \, \text{N} + 490.5 \, \text{N} = 1490.5 \, \text{N} \). The required pressure is then \( P = \frac{F}{A} = \frac{1490.5}{0.01} = 149050 \, \text{Pa} = 1.4905 \, \text{bar} \).
\end{enumerate}

\newpage
\section*{Question 6 [8 points]}
\subsection*{Points Breakdown:}
\begin{itemize}
    \item Calculation of the pressure at the base due to the water: 4 points
    \item Calculation of the pressure of the air confined above the water level: 4 points
\end{itemize}

\subsection*{Solutions:}
\begin{enumerate}
    \item The pressure at the base due to the water is calculated using the formula \( P = \rho \times g \times h \). With \( \rho = 1000 \, \text{kg/m}^3 \), \( g = 9.81 \, \text{m/s}^2 \), and \( h = 30 \, \text{m} \), the pressure is \( P = 1000 \times 9.81 \times 30 = 294300 \, \text{Pa} = 2.943 \, \text{bar} \).
    \item The pressure of the air confined above the water level is determined by subtracting the water pressure from the total pressure: \( P_{\text{air}} = P_{\text{total}} - P_{\text{water}} \). Using the given total pressure of \( 4.15 \, \text{bar} \) and the calculated water pressure of \( 2.943 \, \text{bar} \), the air pressure is \( P_{\text{air}} = 4.15 - 2.943 = 1.207 \, \text{bar} \).
\end{enumerate}

\newpage
\section*{Question 7 [8 points]}
\subsection*{Points Breakdown:}
\begin{itemize}
    \item Calculation of the pressure at the water-oil interface: 4 points
    \item Calculation of the pressure at the bottom of the tank: 4 points
\end{itemize}

\subsection*{Solutions:}
\begin{enumerate}
    \item The pressure at the water-oil interface is calculated as \( P = \rho \times g \times h \). With \( \rho = 890 \, \text{kg/m}^3 \), \( g = 9.81 \, \text{m/s}^2 \), and \( h = 10 \, \text{m} \), the pressure is \( P = 890 \times 9.81 \times 10 = 87309 \, \text{Pa} \) or \( 0.87309 \, \text{bar} \).
    \item The pressure of the water is calculated as \( P = \rho \times g \times h \). With \( \rho = 1000 \, \text{kg/m}^3 \), \( g = 9.81 \, \text{m/s}^2 \), and \( h = 1 \, \text{m} \), the pressure is \( P = 1000 \times 9.81 \times 1 = 9810 \, \text{Pa} \) or \( 0.0981 \, \text{bar} \).
    \item The total pressure at the bottom of the tank is found by adding the pressure due to the water and the pressure due to the oil: \( P_{\text{total}} = P_{\text{water}} + P_{\text{oil}} = 87309 + 9810 = 97119 \, \text{Pa} \) or \( 0.97119 \, \text{bar} \).
\end{enumerate}

\newpage
\section*{Question 8 [6 points]}
\subsection*{Points Breakdown:}
\begin{itemize}
    \item Calculation of the mass of the fluid: 3 points
    \item Calculation of the specific volume: 3 points
\end{itemize}

\subsection*{Solutions:}
\begin{enumerate}
    \item The mass of the fluid is calculated as \( m = \frac{W}{g} \). With \( W = 9810 \, \text{N} \) and \( g = 9.81 \, \text{m/s}^2 \), the mass is \( m = \frac{9810}{9.81} = 1000 \, \text{kg} \).
    \item The specific volume \( \nu \) is calculated as \( \nu = \frac{V}{m} \). With \( V = 150 \, \text{m}^3 \) and \( m = 1000 \, \text{kg} \), the specific volume is \( \nu = \frac{150}{1000} = 0.15 \, \text{m}^3/\text{kg} \).
\end{enumerate}

\newpage
\section*{Question 9 [10 points]}
\subsection*{Points Breakdown:}
\begin{itemize}
    \item Calculation of the mass of the whole system (bottle and contents): 4 points
    \item Calculation of the required force: 6 points
\end{itemize}

\subsection*{Solutions:}
\begin{enumerate}
    \item The mass of the propane is calculated using the molecular weight \( M = 44.1 \, \text{kg/kmol} \) and the amount in kmol, \( n = 1.75 \, \text{kmol} \). Thus, \( m_{\text{propane}} = n \times M = 1.75 \times 44.1 = 77.175 \, \text{kg} \).
    
    \item The total mass of the system, \( m_{\text{total}} \), is the sum of the mass of the propane and the mass of the empty steel bottle \( m_{\text{bottle}} = 12 \, \text{kg} \). Thus, \( m_{\text{total}} = 77.175 + 12 = 89.175 \, \text{kg} \).
    
    \item The force \( F \) required to achieve an acceleration \( a = 3 \, \text{m/s}^2 \) is calculated using \( F = m \times a \). Therefore, \( F = 89.175 \times 3 = 267.525 \, \text{N} \).
\end{enumerate}

\newpage
\section*{Question 10 [10 points]}
\subsection*{Points Breakdown:}
\begin{itemize}
    \item Calculation of the column height: 10 points
\end{itemize}

\subsection*{Solutions:}
\begin{enumerate}
    \item To find the height \( h \) of the column of gasoline when the pressure \( P \) at the bottom is \( 60 \, \text{kPa} \) or \( 60000 \, \text{Pa} \), the density \( \rho \) is \( 715 \, \text{kg/m}^3 \), and the acceleration due to gravity \( g \) is \( 9.81 \, \text{m/s}^2 \), the following formula is used:
    \[
    P = \rho gh \implies h = \frac{P}{\rho g} = \frac{60000}{715 \times 9.81} = 8.55 \, \text{m}
    \]
\end{enumerate}

\newpage
\section*{Question 11 [6 points]}
\subsection*{Points Breakdown:}
\begin{itemize}
    \item Conversion to Fahrenheit: 2 points
    \item Conversion to Kelvin: 2 points
    \item Conversion to Rankine: 2 points
\end{itemize}

\subsection*{Solutions:}
\begin{enumerate}
    \item To convert \( 25^\circ\text{C} \) to Fahrenheit, the following formula is used:
    \[
    T_{\text{F}}^\circ = 1.8 \times 25 + 32 = 77^\circ\text{F}
    \]
    
    \item To convert \( 25^\circ\text{C} \) to Kelvin, the formula is:
    \[
    T_{\text{K}} = 273.15 + 25 = 298.15 \, \text{K}
    \]

    \item To convert \( 298.15 \, \text{K} \) to Rankine, the formula is:
    \[
    T_{\text{R}} = 1.8 \times 298.15 = 536.67 \, \text{R}
    \]
\end{enumerate}

\newpage
\section*{Question 12 [5 points]}
\subsection*{Points Breakdown:}
\begin{itemize}
    \item Conceptual explanation: 2 points
    \item Mathematical proof: 3 points
\end{itemize}

\subsection*{Solutions:}
\subsubsection*{Conceptual Explanation:}
If the pressure in Chamber A and Chamber B are equal, the direction the piston moves will depend on the force acting on either side. This force is governed by the equation \( F = P \times A \).
\subsubsection*{Mathematical Proof:}
Firstly, let us establish the notations and assumptions:
\[
P_A = P_B \quad (\text{Pressure in Chamber A and Chamber B})
\]
\[
A_A \quad (\text{Area of the surface in Chamber A})
\]
\[
A_B \quad (\text{Area of the surface in Chamber B})
\]
The force exerted by the fluid on the piston is:
\[
F = P \times A
\]
In terms of Chamber A and Chamber B, these become:
\[
F_A = P_A \times A_A
\]
\[
F_B = P_B \times A_B
\]
Given that \( P_A = P_B \) and \( A_A < A_B \):
\[
F_A = P_A \times A_A
\]
\[
F_B = P_A \times A_B
\]
Since \( A_A < A_B \), then \( F_A < F_B \).
\subsubsection*{Conclusion:}
The larger force \( F_B \) will dominate, pushing the piston towards Chamber A.

\end{document}